\section{Horizontal and Vertical Canonical Forms}

\begin{definition}[Horizontal canonical form for an MPO]\label{def:horizontal_cf_mpdo}
    \lean{MPOTensor.IsHorizontalCF}
    \notready
    \uses{def:mpo_tensor, def:block_diagonal_tensor, def:injective, def:same_mpv2}
    An MPO tensor is in \emph{horizontal canonical form} if its doubled-index
    MPS tensor generates the same MPV family as a block-diagonal tensor
    $\bigoplus_k \mu_k A_k$ whose blocks are injective, left-canonical, and
    have nonzero weights $\mu_k$.
\end{definition}

\begin{definition}[Vertical canonical form for an MPO]\label{def:vertical_cf_mpdo}
    \lean{MPOTensor.IsVerticalCF}
    \notready
    \uses{def:mpo_tensor, def:bnt, def:same_mpv2, def:block_diagonal_tensor}
    An MPO tensor is in \emph{vertical canonical form} if its diagonal MPS
    tensor $\{M^{ii}\}_{i=0}^{d-1}$ generates the same MPV family as a
    block-diagonal tensor obtained from a basis of normal tensors by repeating
    each block according to a multiplicity and weighting each copy by a
    positive scalar.
    This is the flattened form of the decomposition
    $\bigoplus_\alpha \mu_\alpha \otimes M_\alpha$
    from \cite[Proposition~4.13]{Cirac2016MPDO_arXiv}.
\end{definition}

\begin{remark}[Finite witnesses for block-injective canonical form]%
    \label{rem:word_tuple_span_top}
    \lean{MPSTensor.WordTupleSpanTop}
    \lean{MPSTensor.BlockEntryIndex}
    \lean{MPSTensor.wordEntryFamily}
    \lean{MPSTensor.duplicateScalarBlocks_counterexample}
    \lean{MPSTensor.wordTupleSpanTop_of_wordEntryFamily_linearIndependent}
    \lean{MPSTensor.hasBlockSelectorWords_of_wordTupleSpanTop}
    \lean{MPSTensor.hasBlockSelectorWords_of_wordEntryFamily_linearIndependent}
    \lean{MPSTensor.HasBiCF}
    \lean{MPSTensor.hasBiCF_of_wordTupleSpanTop}
    \lean{MPSTensor.hasBiCF_of_wordEntryFamily_linearIndependent}
    \lean{MPSTensor.HasBlockSelectorWords}
    \lean{MPSTensor.PropBlockInjective}
    \lean{MPSTensor.wordTupleSpanTop_of_propBlockInjective}
    \lean{MPSTensor.hasBiCF_of_propBlockInjective}
    \lean{MPSTensor.duplicateScalarBlocks_not_hasBiCF}
    \lean{MPOTensor.horizontalCFData_of_wordTupleSpanTop}
    \lean{MPOTensor.horizontalCFData_of_wordEntryFamily_linearIndependent}
    \lean{MPOTensor.horizontalCFData_of_propBlockInjective}
    \leanok
    Let $(A_k)$ be a family of block tensors, and fix a length $L$.
    Write $A_k^w$ for the length-$L$ word evaluation of $A_k$ along a word $w$.
    The finite-length spanning hypothesis says that the tuple-valued evaluations
    \[
        w \longmapsto (A_k^w)_k
    \]
    span the full product algebra of block matrices.
    The corresponding finite-length separation property asserts that there
    exists a length $L$ such that if, for every word $w$ of length~$L$,
    \[
        \sum_k \tr(\Delta_k A_k^w) = 0
    \]
    then these pairings vanish on a spanning set for the product algebra, so
    the induced linear functional on the product algebra is zero; by
    nondegeneracy of the product trace pairing, every $\Delta_k$ vanishes.
    In other words, a single finite blocking length already separates the blocks
    by their word evaluations. Combined with injectivity, left-canonicality, and
    nonzero block weights, this finite-length separation yields the horizontal
    canonical-form data of the block-diagonal tensor.

    The entrywise scalar family of block-matrix coefficients encodes the same
    finite-length data. Linear independence of that scalar family implies the
    product-algebra spanning statement, hence a block-injective
    canonical-form witness. The same criterion therefore gives horizontal
    canonical-form data directly.
    The duplicate-scalar counterexample shows that this finite-length witness
    is strictly stronger than blockwise injectivity, left-canonicality, and
    nonzero weights alone: two identical scalar blocks satisfy those hypotheses,
    but the block-injective canonical-form separation condition fails and no
    blocking length makes the entrywise family linearly independent.

    The block-injectivity proposition in
    \cite[lines~340--345]{Cirac2016MPDO_arXiv} gives a concrete sufficient
    criterion for the spanning condition: after a common blocking length each
    block is block-injective, and after a further finite blocking length one
    can form linear combinations of word evaluations which equal the identity
    on one chosen block and vanish on all others. Concatenating an injective
    prefix with this finite block-separation data produces the product-algebra
    spanning condition, hence the horizontal canonical-form data.
\end{remark}

\begin{theorem}[Finite block-separation criterion]
    \label{thm:propblockinj_abstract}
    \lean{MPSTensor.hasBiCF_of_propBlockInjective}
    \leanok
    \uses{rem:word_tuple_span_top}
    If a block family admits the finite-length data described in
    the block-injectivity proposition of
    \cite[lines~340--345]{Cirac2016MPDO_arXiv}, then it is in
    block-injective canonical form.
\end{theorem}

\begin{proof}
    \leanok
    \uses{rem:word_tuple_span_top}
    This is the direct implication from the finite-length data to the
    finite-length trace-separation statement. What remains open in the full
    block-injectivity proposition of
    \cite[lines~340--345]{Cirac2016MPDO_arXiv} is the derivation of this
    finite witness from separated canonical-form / BNT data.
\end{proof}

\begin{lemma}[Blockwise equality from agreement on the MPV family]
    \label{lem:blockwise_insert_eq_of_mpv_agree}
    \lean{MPOTensor.blockwise_insert_eq_of_mpv_agree}
    \leanok
    \uses{def:block_diagonal_tensor, def:injective, thm:propblockinj_abstract}
    Let $\bigoplus_k \mu_k A_k$ be a block-diagonal tensor whose blocks are
    injective, left-canonical, weighted by nonzero scalars, and jointly
    block-injective in the sense that for some blocking length~$L$ the tuple of
    trace pairings against length-$L$ block words is faithful across all blocks
    simultaneously.
    If two physical operators induce the same first-site action on every MPV
    generated by this tensor, then their first-site insertions agree on each
    block $A_k$.
\end{lemma}

\begin{proof}
    \leanok
    \uses{def:block_diagonal_tensor, def:injective, thm:propblockinj_abstract}
    This is Lemma~L of \cite[Appendix~A]{Cirac2016MPDO_arXiv}, using the finite
    criterion from Theorem~\ref{thm:propblockinj_abstract}.
    Expanding the hypothesis over a word of length~$L+1$ beginning with the
    fixed site and folding each block contribution into a trace pairing
    against the corresponding first-site insertion tensor yields, upon taking the
    $Y$-$Z$ difference, a tuple of block matrices whose trace pairings
    against every length-$L$ block word vanish. The block-injective
    canonical-form hypothesis then forces each block difference to be zero, and
    dividing by the nonzero weight $\mu_k^{L+1}$ concludes the blockwise
    equality of the two first-site insertions.
\end{proof}

\begin{lemma}[Matrix product vector of the diagonal tensor]
    \label{lem:mpv_diagonal_tensor}
    \lean{MPOTensor.mpv_diagonalTensor}
    \leanok
    \uses{def:mpo_to_mps}
    For an MPO tensor $M$ and a configuration $\sigma$, the matrix product vector of
    the diagonal tensor $i\mapsto M^{ii}$ equals the matrix product vector of the
    doubled-index tensor $\widehat{M}$ evaluated on the diagonal-paired configuration
    $k\mapsto(\sigma_k,\sigma_k)$:
    \[
        V^{(N)}\bigl(i\mapsto M^{ii}\bigr)(\sigma)
        = V^{(N)}\bigl(\widehat{M}\bigr)\bigl(k\mapsto(\sigma_k,\sigma_k)\bigr).
    \]
\end{lemma}
\begin{proof}\leanok
    At each site value $i$ the diagonal tensor agrees with the doubled-index tensor
    on the diagonal pair,
    \[
        M^{ii} = \widehat{M}^{\,(i,i)}.
    \]
    The matrix product vector is multiplicative over sites, so replacing each factor
    $M^{\sigma_k\sigma_k}$ by $\widehat{M}^{\,(\sigma_k,\sigma_k)}$ reindexes the
    configuration from $\sigma$ to $k\mapsto(\sigma_k,\sigma_k)$, which is the stated
    equality.
\end{proof}

\begin{lemma}[Diagonal tensor as a block decomposition under horizontal form]
    \label{lem:mpv_diagonal_tensor_eq_blocks}
    \lean{MPOTensor.mpv_diagonalTensor_eq_blocks}
    \leanok
    \uses{lem:mpv_diagonal_tensor, def:horizontal_cf_mpdo}
    Suppose the doubled-index tensor $\widehat{M}$ shares its matrix product vector
    family with a block-diagonal assembly $\bigoplus_k \mu_k\, A_k$ of blocks $A_k$ on
    the doubled index, as in a horizontal canonical form. Then the diagonal tensor of
    $M$ generates the same matrix product vector family as the block-diagonal assembly
    $\bigoplus_k \mu_k\, A_k^{\mathrm{diag}}$ on the single-spin index, where
    $A_k^{\mathrm{diag}}(i) = A_k(i,i)$ is the diagonal restriction of each block:
    \[
        V^{(N)}\bigl(i\mapsto M^{ii}\bigr)(\sigma)
        = V^{(N)}\Bigl(\textstyle\bigoplus_k \mu_k\, A_k^{\mathrm{diag}}\Bigr)(\sigma).
    \]
\end{lemma}
\begin{proof}\leanok
    \uses{lem:mpv_diagonal_tensor}
    By Lemma~\ref{lem:mpv_diagonal_tensor} the diagonal matrix product vector equals
    that of $\widehat{M}$ on the diagonal-paired configuration, which by hypothesis
    equals that of the block-diagonal assembly there. Expanding the assembly as the
    weighted sum $\sum_k \mu_k^N\, V^{(N)}(A_k)$ and restricting each block to the
    diagonal pair gives $\sum_k \mu_k^N\, V^{(N)}(A_k^{\mathrm{diag}})$, the matrix
    product vector of the single-spin block-diagonal assembly.
\end{proof}

\begin{lemma}[Diagonal tensor word evaluation]
    \label{lem:eval_word_diagonal_tensor}
    \lean{MPOTensor.evalWord_diagonalTensor}
    \leanok
    \uses{def:mpo_eval_word}
    For any word $w$, the matrix product of the diagonal tensor along $w$ equals the
    operator word evaluation with equal ket and bra words, since both multiply the
    matrices $M^{w_k w_k}$ site by site.
\end{lemma}
\begin{proof}\leanok
    By induction on the word.
\end{proof}

\begin{lemma}[Diagonal tensor evaluates the density-operator diagonal]
    \label{lem:mpv_diagonal_tensor_eq_density}
    \lean{MPOTensor.mpv_diagonalTensor_eq_mpo_diag}
    \leanok
    \uses{def:mpo_family, lem:eval_word_diagonal_tensor}
    For an MPO tensor $M$ and a configuration $\sigma$, the matrix product vector of
    the diagonal tensor equals the diagonal entry of the generated operator,
    \[
        V^{(N)}\bigl(i\mapsto M^{ii}\bigr)(\sigma)
        = \bigl\langle\sigma\,\big|\,\rho^{(N)}(M)\,\big|\,\sigma\bigr\rangle,
    \]
    since both equal $\operatorname{tr}\bigl(M^{\sigma_1\sigma_1}\cdots
        M^{\sigma_N\sigma_N}\bigr)$.
\end{lemma}
\begin{proof}\leanok
    \uses{lem:eval_word_diagonal_tensor}
    Word evaluation of the diagonal tensor multiplies the matrices
    $M^{\sigma_k\sigma_k}$, which is exactly the operator word evaluation with equal
    ket and bra indices; taking the trace gives the diagonal entry.
\end{proof}

\begin{lemma}[Nonnegativity of the diagonal matrix product vector]
    \label{lem:mpv_diagonal_tensor_nonneg}
    \lean{MPOTensor.mpv_diagonalTensor_nonneg}
    \leanok
    \uses{lem:mpv_diagonal_tensor_eq_density, def:mpdo}
    If $\rho^{(N)}(M)$ is positive semidefinite, then
    $V^{(N)}\bigl(i\mapsto M^{ii}\bigr)(\sigma)\ge 0$ for every configuration
    $\sigma$.
\end{lemma}
\begin{proof}\leanok
    By Lemma~\ref{lem:mpv_diagonal_tensor_eq_density} the value is a diagonal entry of $\rho^{(N)}(M)$, and the
    diagonal entries of a positive semidefinite matrix are nonnegative.
\end{proof}

\begin{lemma}[Matrix product vector of the vertical-assembled tensor]
    \label{lem:mpv_vertical_assembled}
    \lean{MPOTensor.mpv_verticalAssembledTensor_eq_sum}
    \leanok
    \uses{thm:mpv_decomposition, def:block_diagonal_tensor}
    Let $g$ be a block count, $\dim,\mult:\{1,\dots,g\}\to\mathbb N$, weights
    $\omega_{\alpha j}\in\mathbb C$ for $j\le\mult_\alpha$, and blocks $A_\alpha$ of
    bond dimension $\dim_\alpha$. The vertical-assembled tensor is the
    block-diagonal assembly over the flattened index $(\alpha, j)$ that repeats
    $A_\alpha$ with weight $\omega_{\alpha j}$, written
    $\bigoplus_\alpha\bigoplus_j \omega_{\alpha j}\, A_\alpha$. Its matrix product
    vector decomposes as
    \[
        V^{(N)}\Bigl(\textstyle\bigoplus_\alpha\bigoplus_j
        \omega_{\alpha j}\, A_\alpha\Bigr)(\sigma)
        = \sum_{\alpha}\sum_{j} (\omega_{\alpha j})^N\, V^{(N)}(A_\alpha)(\sigma).
    \]
\end{lemma}
\begin{proof}\leanok
    \uses{thm:mpv_decomposition}
    The vertical-assembled tensor is the block-diagonal assembly over the single
    flattened index $q$, with flattened weights $\widetilde\omega_q$ and blocks
    $\widetilde A_q$. By the block matrix-product-vector decomposition
    (Theorem~\ref{thm:mpv_decomposition}),
    \[
        V^{(N)}\bigl(\text{assembly}\bigr)(\sigma)
        = \sum_{q} (\widetilde\omega_q)^N\, V^{(N)}(\widetilde A_q)(\sigma).
    \]
    The flattening identifies $q$ with the pair $(\alpha, j)$ via the bijection
    $\bigl(\{1,\dots,g\}\times\{1,\dots,\mult_\alpha\}\bigr)\simeq
        \{1,\dots,\sum_\alpha\mult_\alpha\}$, under which $\widetilde\omega_q=
        \omega_{\alpha j}$ and $\widetilde A_q = A_\alpha$. Reindexing the sum along
    this bijection yields the double sum over $(\alpha, j)$.
\end{proof}

\begin{theorem}[Horizontal canonical form implies vertical canonical form]
    \label{thm:vertical_cf_of_horizontal_cf}
    \notready
    \uses{def:mpdo, def:horizontal_cf_mpdo, def:vertical_cf_mpdo,
        lem:blockwise_insert_eq_of_mpv_agree, lem:mpv_diagonal_tensor,
        lem:mpv_diagonal_tensor_eq_blocks, lem:mpv_diagonal_tensor_nonneg,
        lem:mpv_vertical_assembled}
    Every MPDO in horizontal canonical form is also in vertical canonical form.
\end{theorem}

\begin{proof}
    \uses{def:mpdo, def:horizontal_cf_mpdo, def:vertical_cf_mpdo,
        lem:blockwise_insert_eq_of_mpv_agree, lem:mpv_diagonal_tensor}
    Starting from the horizontal canonical-form decomposition, one uses the MPDO
    positivity condition to compare the first-site action of the doubled-index
    tensor on the diagonal sector. Lemma~\ref{lem:blockwise_insert_eq_of_mpv_agree}
    then shows that the diagonal tensor splits
    along the same normal blocks, with the coefficients rearranged into
    positive scalar weights.
    Flattening those positive diagonal coefficients produces the required
    vertical canonical-form decomposition.
\end{proof}
